\scanpage{14}
\diagram{\begin{tikzcd}[column sep=small]
(9)\quad1\arrow[r,no head]&\widehat\Pi_{g,\nu}\arrow[r,no head]\arrow[d,"1"]&\mathcal M_{g,\nu}\arrow[r,no head]\arrow[d,hook,"?"']&\mathcal N_{g,\nu}\arrow[r]\arrow[d,hook,"?"']&1\\
1\arrow[r,no head]&\widehat\Pi_{g,\nu}\arrow[r]&\unsure{\operatorname{Autl}|\operatorname{Aut}}_{loc}(\widehat\Pi_{g,\nu})\arrow[r]&\operatorname{Autext}_{loc}(\widehat\Pi_{g,\nu})\arrow[r,no head]&1
\end{tikzcd}}

prolongeant (2) et (3). On a donc que l'image $\unsure{\mathcal M'|\mathcal N'}_{g,\nu}$ de $\unsure{\mathcal M|\mathcal N}_{g,\nu}$ dans $\operatorname{Aut}^{\unsure{*}}_{loc}(\unsure{\overline\Pi|\widehat\Pi}_{g,\nu})$ normalise celle $\widehat\Gamma'_{g,\nu}$ de $\widehat\Gamma^+_{g,\nu}$ (la partie ``Teichmüller profini'') et l'on a
\[
(10)\qquad \mathcal M'_{g,\nu}/\Gamma'_{g,\nu}\simeq\text{quotient de }\Gamma_{\overline{\mathbb Q}/\mathbb Q}
\]

\uline{NB} on peut prouver que pour $g$ fixé, $\nu$ grand, c'est un isom. avec $\Gamma_{\overline{\mathbb Q}/\mathbb Q}$ lui-même. Mais je pense que c'est tjrs un iso. \inserted{(conjectures).} C'est le cas pour pour $g=0$ (où le cas critique est $\nu=3$) -- ce qui \deleted{signifie que} \inserted{équivaut aussi que}
\diagram{\begin{tikzcd}
\Gamma_{\overline{\mathbb Q}/\mathbb Q}\arrow[r,hook]&\operatorname{Autext}_{loc}(\widehat\Pi_{0,3})
\end{tikzcd}}
\deleted{$\unsure{M_{0,3}\hookrightarrow\operatorname{Aut}}$}

(\uline{NB}
\diagram{\begin{tikzcd}[column sep=small]
\mathcal N_{0,3}\arrow[r,phantom,"\simeq"]&\widehat\Gamma^+_{0,3}\arrow[d,equal]&{}\times{}&\mathcal N^!_{0,3}\arrow[d,equal]\\
&\Gamma^+_{0,3}=\mathfrak S_3&&\deleted{\unsure{\mathcal N}}\Gamma_{\overline{\mathbb Q}/\mathbb Q}
\end{tikzcd}}
)

et pour $g=1$ (cas critique $\nu=1$)

\mspage{13}
